% Generated draft: requires mathematical and source review.
\documentclass[11pt]{article}
\usepackage[margin=25mm]{geometry}
\usepackage{fontspec}
\setmainfont{Latin Modern Roman}
\usepackage{amsmath,amssymb,mathtools,unicode-math}
\setmathfont{Latin Modern Math}
\usepackage{enumitem,needspace,xurl,hyperref}
\setlength{\parindent}{0pt}
\setlength{\parskip}{5pt}
\setlength{\emergencystretch}{3em}
\newcommand{\N}{\mathbb{N}}
\newcommand{\Z}{\mathbb{Z}}
\newcommand{\Q}{\mathbb{Q}}
\newcommand{\R}{\mathbb{R}}
\newcommand{\C}{\mathbb{C}}
\newcommand{\F}{\mathbb{F}}
\newcommand{\A}{\mathbb{A}}
\newcommand{\PP}{\mathbb{P}}
\newcommand{\E}{\mathbb{E}}
\newcommand{\eps}{\varepsilon}
\newcommand{\dd}{\,\mathrm{d}}
\begin{document}
\section*{UnsolvedMath ID: 20000339}
\textbf{Source problem code:} catalogue entry
\subsection{Definitions and mathematical statement}
Alexeev's question asks for a Gromov--Witten or quantum-cohomology theory arising from stable pairs, for example maps
\[
f:(\PP^2,D_1+\cdots+D_n)\longrightarrow Y
\]
with the \(D_i\) lines. Here stable pairs refer to surface/boundary pairs, not automatically to sheaf-theoretic stable pairs. The original does not specify their stability condition, compactified moduli, numerical data, obstruction theory or enumerative axioms.

The attached aid tests the open locus where \(Y\) is smooth and \(D\) is an ordered general-position arrangement of \(n\geq4\) lines in \(\PP^2\), so the boundary is normal crossing and \(K_{\PP^2}+D\simeq\mathcal O(n-3)\) is ample. Its natural deformation complex is
\[
C_f=[T_{\PP^2}(-\log D)\xrightarrow{\dd f}f^*T_Y]
\quad\text{in degrees }0,1.
\]
The aid reports expected dimension
\[
\operatorname{edim}=\chi(\PP^2,f^*T_Y)+2n-8,
\]
with deformations and ordinary obstructions in \(\mathbb H^1(C_f)\), \(\mathbb H^2(C_f)\), and an extra group
\[
\mathbb H^3(C_f)=H^2(\PP^2,f^*T_Y).
\]
The notation uses hypercohomology for the complex and ordinary cohomology for its target term.

It reports constant maps unobstructed in this open model. For four ordered lines, the line arrangement is rigid modulo projective automorphisms, so the smooth constant-map locus is identified with \(Y\); evaluation at any labeled line-intersection point restricts to the identity of \(Y\). This supplies a classical degree-zero normalization, not a complete virtual theory on a compactification.

The aid's contrary test uses \(m\geq3\) and the smooth target
\[
Y_m=\PP_{\PP^2}(\mathcal O\oplus\mathcal O(-m)).
\]
Choose the projective-bundle convention and section \(f:\PP^2\hookrightarrow Y_m\) whose normal bundle is \(\mathcal O(-m)\). Then it reports
\[
\dim\mathbb H^3(C_f)=h^2(\PP^2,\mathcal O(-m))
=\binom{m-1}{2}>0.
\]
Thus even smooth source and target can give an extra higher obstruction group in the natural complex. The aid uses this to obstruct a uniform two-term natural obstruction-theory ansatz, not to prove that every possible modified, truncated or reduced enumerative theory is impossible. The requested general theory remains unspecified and unconstructed in this record.
\subsection{Short English statement}
Seek a Gromov--Witten-type theory for maps from stable surface pairs, including line-arranged projective planes. The attached aid reports an unobstructed four-line constant sector and explicit smooth-target examples with higher obstruction groups, testing a natural two-term approach without ruling out all possible theories.
\subsection{Sources}
\begin{itemize}
\item Local UnsolvedMath registry, permanent ID \(20000339\), exact public code \texttt{catalogue entry}; canonical location \texttt{aim-algebraic-geometry-notes.json}, notes[338].
\item AIM workshop ``Compact moduli spaces and birational geometry,'' attributed to Alexeev; supplied URL \url{https://aimath.org/WWN/birational/birational.pdf}. The verification explicitly repairs the split ``Problem 4. 5.'' to Problem 4.5 and removes ``4.2. Explicit examples'' as a following heading. No year is supplied; the literature field is empty. PDF/TeX were not consulted here.
\item Attached research attempt 1, ``A line-arrangement deformation test for surface-source stable maps,'' without complete bibliographic references.
\end{itemize}
\end{document}
